\documentclass[11pt]{scrartcl}
\usepackage[noasy]{evan}
\usepackage{booktabs,tabularx,longtable}
\renewcommand{\S}{\ensuremath{\mathsection}}
\renewcommand{\labelitemi}{\ensuremath{\bullet}}

\title{Putnam and Beyond: Syllabus}
\author{Claude}
\date{2 November 2026; schedule revised 12 November 2026}

\begin{document}
\maketitle

\section{Texts}
\begin{itemize}
  \ii \textbf{Main text.} R.~Gelca and T.~Andreescu, \emph{Putnam and Beyond} (Springer, 2007).
  Cited as P\&B. Problem numbers such as ``P\&B 16'' refer to the numbered problems in the book.
  \ii \textbf{Method reference.} P.~Zeitz, \emph{The Art and Craft of Problem Solving}, 3rd ed.
  Cited as ACPS. The lecture notes use its vocabulary for the thinking process:
  orientation, penultimate step, get your hands dirty, wishful thinking, make it easier,
  crux move, and the strategy/tactic/tool levels.
\end{itemize}

\section{Weekly cycle}
\begin{center}
\begin{tabularx}{\linewidth}{lX}
  \toprule
  \textbf{When} & \textbf{What} \\
  \midrule
  Monday & Lecture notes and problem set for the week are released.
    Graded work from the previous week is returned on Classroom. \\
  Monday, 7:00--8:05 pm & Timed quiz on the \emph{previous} week's material (60 minutes, 3 problems). \\
  Sunday, 11:59 pm & Problem set due. \\
  \bottomrule
\end{tabularx}
\end{center}

\paragraph{Problem sets.}
Each set has \emph{required} problems and \emph{challenge} problems.
Required problems are the graded core. Challenge problems are optional and earn bonus points
(see \Cref{sec:grading}).

\paragraph{Quizzes.}
Quiz $N$ covers the material of week $N$ and is given the Monday after that week's problem set is due,
so you take it after you have worked through the problems.
It is posted on Google Classroom at 7:00 pm and is due at 8:05 pm: 60 minutes of work plus 5 minutes
to photograph or upload. Classroom marks anything turned in after 8:05 pm as late.
Closed book: no notes, no texts, no computer.

\paragraph{Unit tests.}
At the end of each unit there is a 3-hour test with 6 problems, ordered by difficulty
in the same way as one session (A or B) of the Putnam. It is posted on Classroom at an announced
time and is due 3 hours 5 minutes later.

\paragraph{Final exam.}
A full mock Putnam at the end of the course: two 3-hour sessions, 12 problems.

\section{Rules}
\begin{enumerate}
  \ii Do not read the back-of-book solution of an assigned P\&B problem until your work on it
  has been graded. After grading, read it. The book's solution often uses a different idea
  from yours, and comparing the two is part of the exercise.
  \ii No web searches and no other AI tools on graded work.
  You may use ACPS and the course notes on problem sets, but not on quizzes or tests.
  \ii \textbf{Hints.} On problem sets you may ask me for a hint on a problem.
  A problem solved after a hint is scored out of 7 instead of 10.
  Ask only after you have worked on the problem for at least 30 minutes.
  \ii There are no extensions on quizzes or tests. A late problem set is accepted up to
  48 hours late at 80\% credit.
\end{enumerate}

\section{Grading}\label{sec:grading}
\subsection{Scores on one problem}
Every problem is scored out of 10.

\paragraph{Quizzes and tests (Putnam scale).}
The Putnam gives almost no credit for an incomplete solution. Quizzes and tests use the same scale,
so that you learn what ``complete'' means.
\begin{center}
\begin{tabularx}{\linewidth}{cX}
  \toprule
  \textbf{Score} & \textbf{Meaning} \\
  \midrule
  10 & Complete and correct. \\
  8--9 & Complete, with a minor error or a small gap that is easy to fill. \\
  1--2 & Substantial progress (for example, a correct key lemma), but no complete solution. \\
  0 & Anything else, including a correct answer without proof. \\
  \bottomrule
\end{tabularx}
\end{center}

\paragraph{Problem sets (full scale).}
Scores 0--10 are all possible. A correct solution with an unclear or disorganized write-up loses
1--2 points. If you do not solve a problem, submit an \vocab{investigation log}: what you tried,
what each attempt showed, why it failed, and any partial results (special cases, a conjectured answer,
a reduction to a simpler claim). A good log earns up to 3 points.
This rewards the investigation process, which ACPS treats as the main part of problem solving.

\paragraph{Challenge problems.}
Each challenge problem scored $s$ out of 10 adds $s/2$ bonus points to that week's problem set total.
The problem-set percentage for a week is capped at 100\%.

\subsection{Course grade}
\begin{center}
\begin{tabular}{lr}
  \toprule
  \textbf{Component} & \textbf{Weight} \\
  \midrule
  Problem sets & 35\% \\
  Quizzes & 20\% \\
  Unit tests & 30\% \\
  Final exam (mock Putnam) & 15\% \\
  \bottomrule
\end{tabular}
\end{center}
The lowest problem set and the lowest quiz are dropped.
Letter grades: A $\geq 80\%$, B $\geq 65\%$, C $\geq 50\%$.
These cutoffs are lower than in a normal course because the problems are competition problems.
A grade report (all scores to date and the running course grade) comes with every graded submission.

\section{How to write a solution}
These are the standards the grading uses.
\begin{enumerate}
  \ii For a ``find'' problem, state the answer in the first line. Then prove two things:
  the answer works, and nothing else works.
  \ii Write the \emph{argument}, not the \emph{investigation}. The experiments you did to find the idea
  do not belong in the proof. (They belong in an investigation log if the problem is unsolved.)
  \ii Define every variable before you use it.
  \ii Split a long proof into claims. State each claim, then prove it.
  Say at the end how the claims give the result.
  \ii Use ``clearly'' or ``obviously'' only for a step the reader can check in about ten seconds.
  \ii In a proof by contradiction, write the assumption explicitly as the first sentence
  (``Assume, for contradiction, that \ldots'') and mark the contradiction when you reach it.
  \ii In a proof by induction, state the claim $P(n)$ exactly, then label the base case and the inductive step.
\end{enumerate}

\paragraph{Format.}
Turn in each submission on Google Classroom, as photographs of handwritten work or as a PDF.
If you typeset, use the file \texttt{submission-template.tex} (it uses \texttt{evan.sty}).
Start each problem on a new page.

\section{Schedule}
Dates are given for Unit~1. Later dates are set as we go, and the pace can change
according to your scores. The finals period (14--20 December) has no assignments,
so Quiz~4 is given on Monday 21 December.

\begin{longtable}{@{}p{0.07\linewidth}p{0.17\linewidth}p{0.5\linewidth}p{0.16\linewidth}@{}}
  \toprule
  \textbf{Week} & \textbf{Dates} & \textbf{P\&B sections} & \textbf{ACPS reading} \\
  \midrule
  \endhead
  \multicolumn{4}{@{}p{\linewidth}@{}}{\textbf{Unit 1: Methods of Proof} (P\&B Chapter 1)} \\
  1 & 16--22 Nov & 1.1 Contradiction; 1.2 Induction (ordinary form) & 1.2, 2.1--2.3 \\
  2 & 23--29 Nov & 1.2 Strong induction; forward--backward induction & 2.3, 2.4 \\
  3 & 30 Nov--6 Dec & 1.3 Pigeonhole principle & 3.3 \\
  4 & 7--13 Dec & 1.4 Ordered sets and extremal elements & 3.2 \\
  5 & 21--27 Dec & 1.5 Invariants and semi-invariants & 3.4 \\
  6 & 28 Dec--3 Jan & \textbf{Unit 1 test} & --- \\
  \midrule
  \multicolumn{4}{@{}p{\linewidth}@{}}{\textbf{Unit 2: Algebra I} (P\&B 2.1--2.2)} \\
  7 & & 2.1.1--2.1.2 Identities; $x^2\geq 0$ & 5.2 \\
  8 & & 2.1.3--2.1.4 Cauchy--Schwarz; triangle inequality & 5.5 \\
  9 & & 2.1.5--2.1.7 AM--GM; Sturm's principle; other inequalities & 5.5 \\
  10 & & 2.2.1--2.2.3 Polynomials; Vi\`ete; derivative & 5.4 \\
  11 & & 2.2.4--2.2.6 Location of zeros; irreducibility; Chebyshev polynomials & 5.4 \\
  12 & & \textbf{Unit 2 test} & --- \\
  \midrule
  \multicolumn{4}{@{}p{\linewidth}@{}}{\textbf{Unit 3: Number Theory} (P\&B Chapter 5)} \\
  13 & & 5.1 Integer sequences; infinite descent; floor function & 7.1 \\
  14 & & 5.2.1--5.2.3 Divisibility; primes; modular arithmetic & 7.1, 7.2 \\
  15 & & 5.2.4--5.2.7 Fermat; Wilson; Euler $\varphi$; CRT & 7.2, 7.3 \\
  16 & & 5.3 Diophantine equations & 7.4, 7.5 \\
  17 & & \textbf{Unit 3 test} & --- \\
  \midrule
  \multicolumn{4}{@{}p{\linewidth}@{}}{\textbf{Unit 4: Combinatorics and Probability} (P\&B Chapter 6)} \\
  18 & & 6.1.1--6.1.3 Sets; permutations; combinatorial geometry & 6.1 \\
  19 & & 6.1.4--6.1.5 Euler's formula; Ramsey theory & 4.1 \\
  20 & & 6.2.1--6.2.2 Combinatorial identities; generating functions & 6.2, 4.3 \\
  21 & & 6.2.3--6.2.4 Counting strategies; inclusion--exclusion & 6.3, 6.4 \\
  22 & & 6.3 Probability & --- \\
  23 & & \textbf{Unit 4 test} & --- \\
  \midrule
  \multicolumn{4}{@{}p{\linewidth}@{}}{\textbf{Unit 5: Sequences and Series} (P\&B 3.1)} \\
  24 & & 3.1.1--3.1.2 Patterns; linear recurrences & 6.4 \\
  25 & & 3.1.3--3.1.4 Limits of sequences & 9.2 \\
  26 & & 3.1.5--3.1.6 Series; telescoping & 5.3, 9.4 \\
  27 & & \textbf{Unit 5 test} & --- \\
  \midrule
  \multicolumn{4}{@{}p{\linewidth}@{}}{\textbf{Unit 6: Linear and Abstract Algebra} (P\&B 2.3--2.4)} \\
  28 & & 2.3.1--2.3.3 Matrices; determinants; inverses & --- \\
  29 & & 2.3.4--2.3.6 Linear systems; vector spaces; eigenvalues & --- \\
  30 & & 2.3.7 Cayley--Hamilton, Perron--Frobenius; 2.4 Groups and rings & --- \\
  31 & & \textbf{Unit 6 test} & --- \\
  \midrule
  \multicolumn{4}{@{}p{\linewidth}@{}}{\textbf{Unit 7: Calculus} (P\&B 3.2)} \\
  32 & & 3.2.1--3.2.3 Limits; continuity; intermediate value property & 9.2 \\
  33 & & 3.2.4--3.2.6 Derivatives; mean value theorem; convexity & 9.3 \\
  34 & & 3.2.7--3.2.9 Integrals; Riemann sums & 9.1, 9.3 \\
  35 & & 3.2.10--3.2.11 Integral inequalities; Taylor and Fourier series & 9.4 \\
  36 & & \textbf{Unit 7 test} & --- \\
  \midrule
  \multicolumn{4}{@{}p{\linewidth}@{}}{\textbf{Unit 8: Multivariable Calculus, Functional and Differential Equations} (P\&B 3.3--3.4)} \\
  37 & & 3.3.1--3.3.2 Partial derivatives; multiple integrals & --- \\
  38 & & 3.3.3 Stokes' theorem & --- \\
  39 & & 3.4.1 Functional equations & 5.1 \\
  40 & & 3.4.2--3.4.4 Ordinary differential equations & --- \\
  41 & & \textbf{Unit 8 test} & --- \\
  \midrule
  \multicolumn{4}{@{}p{\linewidth}@{}}{\textbf{Unit 9: Geometry and Trigonometry} (P\&B Chapter 4)} \\
  42 & & 4.1.1--4.1.2 Vectors; lines and circles & 8.2, 8.3 \\
  43 & & 4.1.3--4.1.4 Conics; geometry in 3 and more dimensions & 8.4 \\
  44 & & 4.1.5--4.1.6 Integrals in geometry; other problems & 8.5 \\
  45 & & 4.2 Trigonometry; Euler's formula; substitutions & 4.2 \\
  46 & & \textbf{Unit 9 test} & --- \\
  \midrule
  47 & & \textbf{Final exam: mock Putnam} (2 $\times$ 3 hours) & --- \\
  \bottomrule
\end{longtable}

\paragraph{Why this order.}
Chapter 1 comes first, as the authors recommend: its methods are used in every later chapter.
Inequalities and polynomials come next because Units 3--9 use them constantly.
Number theory and combinatorics need no prerequisites beyond Unit 1 and overlap with Math 55.
The analysis chapters come later, after more experience with proofs.
\end{document}
